\documentclass[a4paper]{article}

\newcommand{\docTitle}{Robotik}

\usepackage{datetime2}  % for dates
\usepackage[utf8]{inputenc} % UTF-8 input encoding
\usepackage[T1]{fontenc} % fixes \textquotedbl with pdflatex/OT1
\usepackage{textcomp}   % additional text symbols for typewriter/upquote
\usepackage{xcolor}     % for gray color
\usepackage{listings}
\usepackage{marginnote} % for horizontal margin notes
\usepackage{parskip}
\usepackage{amsmath,amssymb,mathtools}
\usepackage{everypage}
\usepackage{fancyhdr}
\usepackage{amsfonts}   % Mengen and natural numbers
\usepackage{enumerate}  % Custom Enumerations
\usepackage[inline]{enumitem}   % List spacing and labels
\usepackage{multicol}   % Multiple Columns
\usepackage{tabularx}

\usepackage{hyperref}   % Links
\hypersetup{
    pdftitle={\docTitle},
    pdfauthor={Moritz Köhler},
    pdfkeywords={DHBW, Hochschule},
    colorlinks=true,
    linkcolor=black,
    filecolor=magenta,
    urlcolor=cyan
}

% -------
% Images
% -------

\usepackage{import}
\usepackage{xifthen}
\usepackage{pdfpages}
\usepackage{tikz}
\usetikzlibrary{arrows.meta, positioning}

\newcommand{\incfig}[1]{%
    \def\svgwidth{\columnwidth}
    \import{./fig/}{#1.pdf_tex}
}

% ---------------------------------
% Vertical Side note (Bottom Left)
% ---------------------------------

\newcommand{\verticaldocinfo}{%
  \rotatebox{90}{\tiny\color{gray}\texttt{\DTMtoday\_\docTitle\_\jobname}}%
}

\AddEverypageHook{
  \put(-40pt,-650pt){\verticaldocinfo}%
}

% ---------------
% Math Shortcuts
% ---------------

% Number sets
\newcommand{\R}{\mathbb{R}}
\newcommand{\N}{\mathbb{N}}
\newcommand{\Z}{\mathbb{Z}}
\newcommand{\Q}{\mathbb{Q}}
\newcommand{\C}{\mathbb{C}}

% Vectors and matrices
\newcommand{\vect}[1]{\mathbf{#1}}
\newcommand{\mat}[1]{\mathbf{#1}}

% Derivatives
\newcommand{\dd}{\mathrm{d}}
\newcommand{\dpartial}[2]{\frac{\partial #1}{\partial #2}}
\newcommand{\ddt}{\frac{\dd}{\dd t}}

% Norms and absolute values
\newcommand{\norm}[1]{\left\lVert #1 \right\rVert}
\newcommand{\abs}[1]{\left\lvert #1 \right\rvert}

% Probability & Statistics
\newcommand{\E}{\mathbb{E}}
\newcommand{\Var}{\mathrm{Var}}
\newcommand{\Prob}{\mathbb{P}}

% Fields and Algebras
\newcommand{\K}{\mathbb{K}}
\newcommand{\F}{\mathbb{F}}

% Set Operations
\newcommand{\compl}[1]{#1^{\mathrm{c}}}

% Matrix and Vector Operations
\newcommand{\T}{^{\top}}
\newcommand{\inv}{^{-1}}
\newcommand{\conj}[1]{\overline{#1}}

% -------------------
% Main lecture macro 
% -------------------

\newcommand{\lecture}[1]{%
  \protect\marginnote{\protect\tiny\protect\color{gray}#1}
}

% -----------------------------
% Command for box with heading
% -----------------------------

\fboxsep2mm
\newcommand{\know}[2]{
    \noindent\fbox{
        \parbox{\textwidth-21pt}{

            \textbf{#1:}
            \vspace{0.125cm}

            #2
        }
    }
}

% -------------------
% Listings Languages
% -------------------

% Shared defaults for code listings
\lstset{
    basicstyle=\ttfamily\small,
    keywordstyle=\color{blue}\bfseries,
    commentstyle=\color{gray},
    stringstyle=\color{teal},
    numberstyle=\tiny\color{gray},
    numbers=left,
    stepnumber=1,
    numbersep=8pt,
    showstringspaces=false,
    frame=single,
    breaklines=true,
    tabsize=4,
    keepspaces=false,
    columns=flexible,
    upquote=true,
    extendedchars=false % disable extended char handling to avoid UTF-8 conflicts
}

% SQL syntax highlighting
\lstdefinelanguage{SQL}{
    morekeywords={SELECT,FROM,WHERE,INSERT,INTO,VALUES,UPDATE,SET,DELETE,CREATE,TABLE,PRIMARY,KEY,NOT,NULL,AND,OR,JOIN,LEFT,RIGHT,INNER,ON,AS,ORDER,BY,GROUP,HAVING,LIMIT},
    sensitive=false,
    morecomment=[l]{--},
    morestring=[b]'
}

% JSON syntax highlighting
\lstdefinelanguage{json}{
    morestring=[b]",
    stringstyle=\color{teal},
    comment=[l]{//},
    morecomment=[s]{/*}{*/},
        morekeywords={true,false,null},
    literate=
     *{0}{{{\color{blue}0}}}{1}
      {1}{{{\color{blue}1}}}{1}
      {2}{{{\color{blue}2}}}{1}
      {3}{{{\color{blue}3}}}{1}
      {4}{{{\color{blue}4}}}{1}
      {5}{{{\color{blue}5}}}{1}
      {6}{{{\color{blue}6}}}{1}
      {7}{{{\color{blue}7}}}{1}
      {8}{{{\color{blue}8}}}{1}
      {9}{{{\color{blue}9}}}{1}
            {:}{{{\color{black}:}}}{1}
            {,}{{{\color{black},}}}{1},
    showstringspaces=false
}

% Language specific styles
\lstdefinestyle{javastyle}{
    language=Java,
    basicstyle=\ttfamily\small,
    keywordstyle=\color{blue}\bfseries,
    commentstyle=\color{gray},
    stringstyle=\color{teal},
    numberstyle=\tiny\color{gray},
    numbers=left,
    stepnumber=1,
    numbersep=8pt,
    showstringspaces=false,
    frame=single,
    breaklines=true,
    tabsize=4,
    keepspaces=true,
    columns=fullflexible,
    upquote=true
}

\lstdefinestyle{sqlstyle}{
    language=SQL,
    basicstyle=\ttfamily\small,
    keywordstyle=\color{blue}\bfseries,
    commentstyle=\color{gray},
    stringstyle=\color{teal},
    numberstyle=\tiny\color{gray},
    numbers=left,
    stepnumber=1,
    numbersep=8pt,
    showstringspaces=false,
    frame=single,
    breaklines=true,
    tabsize=4,
    keepspaces=true,
    columns=fullflexible,
    upquote=true
}

\lstdefinestyle{pythonstyle}{
    language=Python,
    basicstyle=\ttfamily\small,
    keywordstyle=\color{blue}\bfseries,
    commentstyle=\color{gray},
    stringstyle=\color{teal},
    numberstyle=\tiny\color{gray},
    numbers=left,
    stepnumber=1,
    numbersep=8pt,
    showstringspaces=false,
    frame=single,
    breaklines=true,
    tabsize=4,
    keepspaces=true,
    columns=fullflexible,
    upquote=true
}

\lstdefinestyle{cstyle}{
    language=C,
    basicstyle=\ttfamily\small,
    keywordstyle=\color{blue}\bfseries,
    commentstyle=\color{gray},
    stringstyle=\color{teal},
    numberstyle=\tiny\color{gray},
    numbers=left,
    stepnumber=1,
    numbersep=8pt,
    showstringspaces=false,
    frame=single,
    breaklines=true,
    tabsize=4,
    keepspaces=true,
    columns=fullflexible,
    upquote=true
}

\lstdefinestyle{bashstyle}{
    language=bash,
    basicstyle=\ttfamily\small,
    keywordstyle=\color{blue}\bfseries,
    commentstyle=\color{gray},
    stringstyle=\color{teal},
    numberstyle=\tiny\color{gray},
    numbers=left,
    stepnumber=1,
    numbersep=8pt,
    showstringspaces=false,
    frame=single,
    breaklines=true,
    tabsize=4,
    keepspaces=true,
    columns=fullflexible,
    upquote=true
}

\lstdefinestyle{jsonstyle}{
    language=json,
    basicstyle=\ttfamily\small,
    keywordstyle=\color{blue}\bfseries,
    commentstyle=\color{gray},
    stringstyle=\color{teal},
    numberstyle=\tiny\color{gray},
    numbers=left,
    stepnumber=1,
    numbersep=8pt,
    showstringspaces=false,
    frame=single,
    breaklines=true,
    tabsize=4,
    keepspaces=true,
    columns=fullflexible,
    upquote=true
}

% Default listing style
\lstset{language={}}

\title{\docTitle}
\author{Moritz}
\date{\today}

\begin{document}
\maketitle
\tableofcontents

\lecture{2026-28-09}

\section{Start}

Wir machen viel Übung und Vorlesung und als Labor etwas mit einem Mehrachsenroboter arbeiten.

\subsection{Prüfungsleistung}

Prüfung enthält teile von QuizAcademy Fragen

MATLAB Onramp macht 15\% der Prüfungsleistung aus

Wir sollten R2026a installieren. 

Idealerweise alle Toolboxen, mindestens Toolboxen:

\begin{itemize}
    \item Stateflow
    \item Robotics
    \item Simscape
\end{itemize}

Übungen sind unbewertet. Die zweiter Übung ist Voraussetzung für das Labor. Wir müssen das Labor mit der Übung vorbereiten.

Labore geben noch 35\% geben. Zu zweit einen Roboter. Labore werden mittels Demos geprüft. Also einfach melden und Unterschriften Sammeln. Es gibt keine Laborberichte.

Unser Roboter ist der Braccio mit 6 Freiheitsgraden (Eine ist der Greifer)

Dann noch eine Prüfung 50\% mit 30min

\begin{itemize}
    \item 10 multiple Choice Fragen
    \item Kleine Rechnung
\end{itemize}

Wir dürfen Formelsammlung drei Seiten Doppelt mitnehmen und Taschenrechner.

\subsection{Quellen}

P. Corke, W. Jachimczyk, R. Pillat 2023: Robotics, Vision and Control

B. Siciliano, et al. 2003: Robotics - Modelling, Planning and Control

weitere auf den Folien.

\subsection{Übersicht}

\begin{enumerate}
    \item Einführung
    \item Sensorik, Aktorik und Steuerungstechnik
    \item Modellierung und Simulation
    \item Planung, Vorsteuerung und Regelung
    \item Mobile Roboter
    \item Computer Vision
    \item Robotik und Künstliche Intelligenz
\end{enumerate}

\subsection{Termine}

Labor entweder 2.11 Gruppe A oder 9.11 Gruppe B

Am 7.12 um 13:00 ist die Prüfung

\section{Einführung}

Drei Robotergesetze von Asimov. Die Anwendbarkeit ist fragwürdig

\begin{enumerate}
    \item Menschen wissentlich und durch Untätigkeit nicht verletzten
    \item Muss Befehlen von Menschen gehorchen außer es verstößt 1.
    \item Muss eine eigne Existenz schützen außer 1. und 2.
\end{enumerate}

Im Militär gelten diese Gesetze nicht.

\know{Definition}{"A robot is a goal-oriented machine that can sense, plan, and act." - P. Corke, et al.}

\subsection{Zukunft der Robotik}

Entwicklungstrends:

\begin{enumerate}
    \item Physical AI
    \item Agentic AI
    \item Humanoide werden marktreif
    \item IT/OT-Konvergenz
    \item Cobots
    \item Cyber-Security
\end{enumerate}

\subsection{Grundlagen der Robotik}

Moravec'sches Paradoxon.

Klassifikation nach Anwendung:

\begin{itemize}
    \item Industrieroboter
    \item Cobots
    \item Mobile Roboter (Auto, Drohne, ...)
    \item Service-Roboter
    \item Nano-Roboter
    \item Humanoide Roboter
    \item ...
\end{itemize}

Klassifikation nach Kinematik

\begin{itemize}
    \item Robot mit kartesischen Koordination (Fräse)
    \item Roboter mit zylindrischen Koordination
    \item Knickarmroboter
    \item Parallelogram-Knickarmroboter
    \item SCARA-Roboter
    \item Parallel-Kinematik (Delta-Roboter)
\end{itemize}

\subsection{Herausforderungen der Robotik}

Diskussion: WElche Herausforderungen können im beschriebenen Szenario auftreten?

Es werden kleine und große Boxen auf einem Förderband angeliefert. Bei Erkennung einer Box im Entladebereich mit Hilfe einer Lichtschranke wird das Förderband gestoppt. Dann soll ein Mehrgelenkroboter mit 5 Freiheitsgraden die Boxen auf dem jeweils passenden Förderband zum Einzeltransport positionieren.

\begin{enumerate}
    \item Benennen Sie die 5 Freiheitsgrade des Roboters.
    
    Drehung (Base), Schulter (Shoulder), Ellbogen (Elbow), Handgelenk (Wrist) und Greifen (Gripper).
    \item Welche Herausforderungen bestehen bei der Lösung der gegebenen Aufgabenstellung?
    
    \begin{itemize}
        \item Unterscheidung der Boxen anhand ihrer Größe.
        \item Greifen der Boxen ohne Beschädigung
        \item Inverse Kinematics (Bewegung des Arms)
        \item Bahnplanung/ Trajektorenplanung.
        \item Ist das Förderband wo die Box hin soll frei?
        \item (Was passiert mit Boxen die nicht auf die weiteren Förderbände gehören?)
        \item Betriebsmodi, Notaus? (Lösung Zustandsautomat)
    \end{itemize}
    \item Wie kann die zugrundeliegende Automatisierungaufgabe alternativ gelöst werden?
    
    Ein Schiebearm (Sortierweiche), der die Boxen auf ein Linkes oder Rechtes Förderband wegschiebt (Wie Pfandautomat)

    Oder auch Mechanisch durch die größenunterschied der Boxen

    Problem Vorverlagen, das es gar nicht erst gemischt ist.
\end{enumerate}

\end{document}